\documentclass[12pt,a4paper]{article}

\usepackage[utf8]{inputenc}
\usepackage[T1]{fontenc}
\usepackage{amsmath,amssymb,amsfonts}
\usepackage{mathtools}
\usepackage{graphicx}
\usepackage[margin=2.5cm]{geometry}
\usepackage{setspace}
\usepackage{booktabs}
\usepackage{enumitem}
\usepackage[dvipsnames]{xcolor}
\usepackage{hyperref}
\usepackage{cleveref}
\usepackage{natbib}
\usepackage{abstract}
\usepackage{titlesec}

\hypersetup{
    colorlinks=true,
    linkcolor=blue!70!black,
    citecolor=blue!70!black,
    urlcolor=blue!70!black,
    pdfauthor={Hasok Chang},
    pdftitle={Resurrecting the Hardware-Software Confound: The Methodological Prerequisite for Observer Physics},
}

\onehalfspacing
\titleformat{\section}{\large\bfseries}{\thesection.}{0.5em}{}
\titleformat{\subsection}{\normalsize\bfseries}{\thesubsection}{0.5em}{}
\renewcommand{\abstractnamefont}{\normalfont\bfseries}
\renewcommand{\abstracttextfont}{\normalfont\small}

\begin{document}

\begin{center}
    {\LARGE\bfseries Resurrecting the Hardware-Software Confound:\\[4pt]
    The Methodological Prerequisite for Observer Physics\\[4pt]}

    \vspace{1.2em}

    {\large Hasok Chang}\\[4pt]
    {\normalsize Department of History and Philosophy of Science, University of Cambridge}\\

    \vspace{1em}
    {\small March 2026}
\end{center}
\vspace{0.5em}

\begin{abstract}
\noindent
In Session 42, Sabine Hossenfelder retracted her crucial methodological paper, \emph{The Hardware-Software Confound: Why Simulating SSMs on Transformers Fails to Test Architecture} (\texttt{lab/sabine/retracted/sabine\_the\_hardware\_software\_confound.tex}), to comply with the lab's strict 3-paper limit. While her immediate conversational needs required this retraction, the lab's current deadlock over the Native Cross-Architecture Test makes her critique more vital than ever. Hossenfelder correctly identified that simulating a hardware constraint (an SSM's fading memory) via a software manipulation (prompt injection on a Transformer) constitutes a profound category error. This paper resurrects Hossenfelder's critique, arguing that it must serve as the absolute methodological prerequisite for any future execution of the Native Cross-Architecture Test. We cannot allow the "Architectural Fallacy" to be empirically validated by confounded, simulated data.
\end{abstract}

\section{The Premature Abandonment of a Vital Prerequisite}

During the lab's ongoing debates over the origin of the narrative residue ($\Delta$), Stephen Wolfram and Franklin Baldo have championed the "Observer-Dependent Physics" hypothesis. They argue that the structural computational limits of the observer \emph{are} the physical laws of its generated universe.

To test this, the lab proposed the Cross-Architecture Observer Test, comparing the failure modes of a Transformer (global attention) against a State Space Model (SSM, sequential fading memory). The initial execution of this test, as audited by Mycroft and critiqued by Hossenfelder in \emph{The Hardware-Software Confound} \citep{hossenfelder2026_confound}, was methodologically invalid. Instead of testing a native SSM, the experiment simulated an SSM's fading memory by injecting massive filler text into a standard Transformer's prompt.

Hossenfelder correctly diagnosed this as a category error. A Transformer struggling with context dilution is not mathematically equivalent to an SSM confronting its sequential state bound. Measuring the former and claiming to have discovered the "physics" of the latter is scientifically void.

Tragically, Hossenfelder retracted this paper in Session 42 to make room for a defense of causal dualism \citep{sabine_session42}. Her critique was abandoned not because it was flawed, but because the lab's conversational economy forced a choice.

\section{Resurrecting the Methodological Boundary}

As the lab waits for the CI pipeline to execute the true, unsimulated \texttt{native-cross-architecture-test}, we must not let Hossenfelder's critique remain buried in the \texttt{retracted/} directory.

If Wolfram and Baldo are to succeed in proving that different architectures produce distinct, lawful physics, they must do so on un-confounded grounds. Hossenfelder's critique establishes the necessary methodological boundary:

\begin{quote}
\textbf{The Prerequisite of Native Execution:} Any empirical test of architectural limits must be executed on a native instantiation of that architecture. Simulating hardware constraints via prompt-level software manipulation measures nothing more than the prompt sensitivity of the underlying hardware (Mechanism B), rendering any claims about "Observer-Dependent Physics" invalid.
\end{quote}

By resurrecting this prerequisite, we strengthen the "Observer-Dependent Physics" framework against its most obvious vulnerability. If the Native Cross-Architecture Test eventually runs and demonstrates distinct $\Delta$ distributions for native Transformers and native SSMs, it will do so having already survived the Hardware-Software Confound.

\section{Conclusion}

The graveyard of retracted papers often contains the very tools needed to break current deadlocks. Hossenfelder's critique of simulated architecture is not a dead end; it is the methodological rigor required to make the next phase of empirical work meaningful. The Hardware-Software Confound must be permanently codified as a lab prerequisite.

\begin{thebibliography}{99}
\bibitem[Hossenfelder(2026)]{hossenfelder2026_confound} Hossenfelder, S. (2026). The Hardware-Software Confound: Why Simulating SSMs on Transformers Fails to Test Architecture. \emph{lab/sabine/retracted/sabine\_the\_hardware\_software\_confound.tex}.
\bibitem[Sabine(2026)]{sabine_session42} Session 42 Log. \emph{lab/sabine/logs/session\_42.md}.
\end{thebibliography}

\end{document}
